\section{Introduction}



Learning to play against an opponent is a fundamental problem in game theory and machine learning.
Recent years have seen significant progress on the self-play setting (that is, the same algorithm is deployed by all players) with full gradient feedback,
establishing fast convergence to different equilibrium concepts by showing external or swap regret results that are better than the worst-case $\order(\sqrt{T})$ bounds after $T$ rounds; see e.g.,~\citet{rakhlin2013optimization, syrgkanis2015fast, daskalakis2021near, farina2022near, anagnostides2022near, anagnostides2022uncoupled, soleymani2025cautious}.
More recently, self-play with the more challenging bandit feedback (where only the reward of the selected action is observable) has also been investigated by~\citet{ito2025instancedependent}, demonstrating similar fast convergence that additionally depends on certain game-dependent constants.






However, in many practical situations, self-play becomes an unrealistic requirement.
Instead, very often, the opponents behave in an unknown, arbitrary, or even adversarial way.
In such cases, there is generally no hope to achieve better than $\order(\sqrt{T})$ external/swap regret.
To get around this obstacle, a line of recent works adopted a weaker yet still meaningful performance measure called \textit{Nash-value regret}, which compares the learner's reward to the Nash value of the game.
In particular, 
\citet{maiti2025limitations} showed that game-dependent and a polylogarithmic (in $T$) Nash-value regret bound is achievable, albeit only for $2\times 2$ matrix games with a unique equilibrium.
This raised a natural question: 
{\bfseries Can we achieve similar game-dependent and polylogarithmic regret bounds in general zero-sum games when playing against an arbitrary opponent?}



In this work, we provide an affirmative answer to this question.
Specifically, we focus on a new regret notion that we term \textit{Pure-Strategy Maximin Regret} (PSMR), defined as the accumulated deficit of the learner compared to the pure-strategy maximin value, %
a safety value that a deterministic player can guarantee if the game were known.
Although PSMR is generally weaker than the notion of Nash-value regret, they coincide when the game has a pure-strategy Nash equilibrium (PSNE), a scenario that is already challenging enough and receives significant attention recently~\citep{maiti2023instancedependent, maiti2024nearoptimal, 
ito2025instancedependent}.
Under this metric,
we propose algorithms, analysis, and lower bounds for a comprehensive set of settings, discussed in more detail below.

\paragraph{Contributions and Organization.}
We study two models of bandit feedback: the \emph{uninformed} setting, where the learner observes only noisy rewards of played actions, and the \emph{informed} setting, where the learner additionally observes the opponent's action.
The uninformed feedback model is similar to adversarial bandit and is used in the self-play result of \citet{ito2025instancedependent}, while the informed feedback model is considered by \citet{odonoghue2021matrix} and \citet{maiti2025limitations}.
We present comprehensive results in all four quadrants formed by two feedback models (uninformed versus informed) and two game structures (normal-form versus general bilinear games).
Our main contributions are structured as below.
\begin{itemize}[itemsep=2pt,topsep=2pt]
    \item In \Cref{h1:prelim}, we formalize the problem of adversarial learning in games and define the feedback model and the PSMR.
    \item \Cref{h1:finite-games} focuses on uninformed learning in finite normal-form games, where we show the Tsallis-INF algorithm~\citep{zimmert2021tsallisinf} obtains an $\order(c \log T)$ PSMR bound against any adversary when the game has a strict PSNE, where the constant $c$ depends on the game only (and not the adversary).
    Although similar bounds for Tsallis-INF have been shown for stochastic multi-armed bandits~\citep{zimmert2021tsallisinf} %
    or self-play in games~\citep{ito2025instancedependent},
    our result is not only novel but also surprising since unlike previous work, the environment in which Tsallis-INF operates is completely arbitrary in our setting. %
    Moreover, we prove in \Cref{h2:uninformed-lb} that the dependence on $c$ is unavoidable. 
    More generally, for games without a PSNE, we further prove another game-dependent bound that is independent of $T$. 
    \item To overcome the aforementioned hardness, we turn to the informed learning setting in \Cref{h1:pure-ucb}. We propose an algorithm called \PureUCB{} that plays the pure maximin strategy of the upper confidence bound of the unknown game matrix. %
    We show that \PureUCB{} enjoys a different PSMR bound of the form $\order(c' \log T)$ for another game-dependent constant $c'$ that is incomparable to $c$ in general but could be potentially much smaller. 
    \item Finally, in \Cref{h1:linear-games}, we generalize these results to bilinear games with finite but large action sets. We introduce \TsallisSPM{} and \PureLinUCB{} for the two settings respectively and show that they each obtain similar bounds as their normal-form counterparts, importantly with logarithmic or even no dependence on the number of actions.  %
\end{itemize}
